% PDF 85--91；原书第 55--61 页。
\subsection{架构}
\refcite{ch02-vaswani}{Vaswani 等人（2017）}最初提出的 Transformer，由堆叠的编码器层和解码器层组成。预训练基于 Transformer 的语言模型时，可以堆叠编码器层、解码器层，也可以两者并用，由此得到多种架构，如图~\ref{fig:ch02-patterns} 所示。

\subsubsection{编码器—解码器}
编码器—解码器架构分为两部分。编码器处理输入序列，将输入变成一个包含所学输入信息的连续表示。解码器再用这个表示生成输出序列。该架构适用于机器翻译、文本摘要等序列到序列任务。例如，在机器翻译中，编码器处理源语言句子，解码器生成目标语言译文。该架构的注意力机制使模型在生成输出时能够关注输入序列的不同部分，动态计算所需的上下文。

\subsubsection{因果解码器}
因果解码器架构面向自回归任务，模型在这类任务中逐个生成词元。它采用单向注意力机制，即生成时每个词元只能关注前面的词元及其自身。对于需要生成连贯且符合上下文的文本生成任务，这一点尤其有用。例如，补全文本时，模型根据前面的词元预测下一个词元，使生成的文本连贯并与上下文相关。

\begin{figure}[!htbp]
\centering
% 重绘自原书第 56 页的图 2.8；轴标签是原示例中的六个英文词元。
\begingroup
\colorlet{prefixcell}{blue!24}
\colorlet{crosscell}{green!32!gray!13}
\colorlet{targetcell}{yellow!37!orange!12}
\colorlet{maskedcell}{blue!12!gray!37}
\newcommand{\attentiontiles}[2]{%
  \node[font=\footnotesize\sffamily\bfseries] at (1.32,.49) {#2};
  \foreach \row/\ca/\cb/\cc/\cd/\ce/\cf in {#1} {
    \foreach \col/\tone in {0/\ca,1/\cb,2/\cc,3/\cd,4/\ce,5/\cf} {
      \node[tile,fill=\tone] at ({.53*\col},{-.53*\row}) {};
    }
  }
  \foreach \word [count=\idx] in {A,Survey,of,Large,Language,Models} {
    \node[anchor=east] at (-.24,{-.53*(\idx-1)}) {\word};
    \node[rotate=90,anchor=east] at ({.53*(\idx-1)},-2.91) {\word};
  }
}
\newcommand{\decoderbrackets}{%
  \draw[group] (-1.60,.17) -- (-1.73,.17) -- (-1.73,-2.82) -- (-1.60,-2.82);
  \node[rotate=90] at (-1.97,-1.32) {解码器};
  \draw[group] (-.19,-3.94) -- (-.19,-4.08) -- (2.83,-4.08) -- (2.83,-3.94);
  \node at (1.32,-4.29) {解码器};
}
\begin{tikzpicture}[
  font=\scriptsize\sffamily,
  tile/.style={draw=black!56,minimum size=4.1mm,inner sep=0pt,rounded corners=1pt,line width=.48pt},
  group/.style={draw=black!49,line width=.52pt,rounded corners=1pt}]
\begin{scope}[shift={(2.74,5.13)}]
\attentiontiles{
  0/prefixcell/maskedcell/maskedcell/maskedcell/maskedcell/maskedcell,
  1/prefixcell/prefixcell/maskedcell/maskedcell/maskedcell/maskedcell,
  2/prefixcell/prefixcell/prefixcell/maskedcell/maskedcell/maskedcell,
  3/crosscell/crosscell/crosscell/targetcell/maskedcell/maskedcell,
  4/crosscell/crosscell/crosscell/targetcell/targetcell/maskedcell,
  5/crosscell/crosscell/crosscell/targetcell/targetcell/targetcell%
}{因果解码器}
\decoderbrackets
\end{scope}
\begin{scope}
\attentiontiles{
  0/prefixcell/prefixcell/prefixcell/maskedcell/maskedcell/maskedcell,
  1/prefixcell/prefixcell/prefixcell/maskedcell/maskedcell/maskedcell,
  2/prefixcell/prefixcell/prefixcell/maskedcell/maskedcell/maskedcell,
  3/crosscell/crosscell/crosscell/targetcell/maskedcell/maskedcell,
  4/crosscell/crosscell/crosscell/targetcell/targetcell/maskedcell,
  5/crosscell/crosscell/crosscell/targetcell/targetcell/targetcell%
}{前缀解码器}
\decoderbrackets
\end{scope}
\begin{scope}[xshift=5.50cm]
\attentiontiles{
  0/prefixcell/prefixcell/prefixcell/maskedcell/maskedcell/maskedcell,
  1/prefixcell/prefixcell/prefixcell/maskedcell/maskedcell/maskedcell,
  2/prefixcell/prefixcell/prefixcell/maskedcell/maskedcell/maskedcell,
  3/crosscell/crosscell/crosscell/targetcell/maskedcell/maskedcell,
  4/crosscell/crosscell/crosscell/targetcell/targetcell/maskedcell,
  5/crosscell/crosscell/crosscell/targetcell/targetcell/targetcell%
}{编码器—解码器}
\draw[group] (-1.60,.17) -- (-1.73,.17) -- (-1.73,-1.26) -- (-1.60,-1.26);
\draw[group] (-1.60,-1.38) -- (-1.73,-1.38) -- (-1.73,-2.82) -- (-1.60,-2.82);
\node[rotate=90] at (-1.97,-.55) {编码器};
\node[rotate=90] at (-1.97,-2.10) {解码器};
\draw[group] (-.19,-3.94) -- (-.19,-4.08) -- (1.26,-4.08) -- (1.26,-3.94);
\draw[group] (1.38,-3.94) -- (1.38,-4.08) -- (2.83,-4.08) -- (2.83,-3.94);
\node at (.53,-4.29) {编码器};
\node at (2.12,-4.29) {解码器};
\end{scope}
\end{tikzpicture}
\endgroup

\caption{三种主要架构的注意力模式。蓝色、绿色、黄色和灰色圆角方块分别表示前缀词元内部的注意力、前缀词元与目标词元之间的注意力、目标词元内部的注意力以及被掩码遮盖的注意力。}
\label{fig:ch02-patterns}
\end{figure}

\subsubsection{前缀解码器}
前缀解码器是因果解码器的变体。模型可以在一段作为前缀的词元内部双向计算注意力，对其余词元仍保持单向注意力。采用这种混合机制，模型每生成一个词元时都可以利用更广的上下文，因而适合需要理解序列中前后词元的任务。例如，在对话系统中，生成下一个词元时，模型可以关注对话历史和已经生成的部分回复。

\subsubsection{编码器}
编码器能够高效处理并理解输入序列中的上下文信息，适合某些自然语言处理任务。各编码器层都会为输入序列生成能够反映上下文的表示。最后一个编码器层的最终输出用作上下文表示，为多种下游任务提供输入。编码器架构特别适合需要深入理解词元上下文、无须生成序列的任务，例如分类任务。

\subsubsection{混合专家}
混合专家（Mixture of Experts，MoE）架构是 Transformer 的一种变体。它用混合专家层替换标准前馈模块，如图~\ref{fig:ch02-moe} 所示。每个混合专家层包含多个并行单元，称为“专家”，每个专家都有各自的参数。路由器（router）根据专家的能力，把输入词元交给相应专家。专家是注意力模块后的前馈层，彼此独立地处理词元。传统模型扩大容量会增加计算成本，而混合专家架构每次只激活少量专家。这种稀疏激活让模型规模能够继续增大，而计算需求无需等比例增加，因而仍能高效运行。

\sourcefigure{fig02_09.pdf}{Transformer 架构的混合专家变体。左侧路由器（Router）将输入 $x_1$ 送至 FFN 2（$p=0.65$），将输入 $x_2$ 送至 FFN 1（$p=0.80$）。右侧标出混合专家层（MoE layer）、层归一化（Layer Norm）和掩码自注意力（Masked Self-Attention）的位置。}{fig:ch02-moe}

\subsection{大语言模型预训练的技巧与策略}
本节讨论训练大语言模型时常用的配置、方法和策略。

\subsubsection{训练方法}
\begin{itemize}
\item \textbf{学习率。}大多数大语言模型预训练时采用相近的学习率计划，分为预热（warm-up）和衰减阶段。先用约 0.1\% 至 0.5\% 的训练步数逐渐增大学习率，学习率通常在 $5\times 10^{-5}$ 至 $1\times 10^{-4}$ 的范围。此后采用余弦衰减（cosine decay）策略逐步减小学习率。
\item \textbf{批量大小。}语言模型预训练常用大批量，如 2\,048 个样例或 400 万个词元，以改善训练的稳定性和效率。GPT-3 和 PaLM 等模型会在训练中动态调整批量大小。例如，GPT-3 的批量大小从 3.2 万词元增至 320 万词元。研究表明，适应性调整批量大小能有效使大语言模型训练更平稳。
\item \textbf{优化器。}训练 GPT-3 等大语言模型时，常使用 Adam 和 AdamW。这些优化器根据梯度估计调整更新，典型超参数（hyperparameter）为 $\beta_1=0.9$、$\beta_2=0.95$ 和 $\epsilon=10^{-8}$。PaLM、T5 等模型还采用 Adafactor 优化器，这是内存用量较少的 Adam 变体。它的超参数 $\beta_1=0.9$，$\beta_2$ 则随训练步数调整。
\end{itemize}

\subsubsection{解码策略}
\paragraph*{贪心搜索}
自回归解码（autoregressive decoding）是纯解码器架构所用的技术之一。贪心搜索（greedy search）是最常用的解码方法之一。它在每一步根据此前生成的词元，预测概率最大的词元，数学形式如下。
\[
x_i=\arg\max_x P(x\mid x_{<i}),
\]
其中 $x_i$ 是第 $i$ 步预测的词元，即给定上下文 $x_{<i}$ 时概率最大的词元。以未完成的句子“The sky is so”为例，“blue”很可能恰当地补全这个句子，贪心搜索可能把它预测为下一个词元。在机器翻译和文本摘要等文本生成任务中，输入与预期输出之间有较强的依赖，这种办法很有效。

在输出需要紧扣输入的情形中，贪心搜索利用概率和上下文给出可靠的结果。这种解码策略也可以用于编码器—解码器和前缀解码器模型。研究者提出了许多改进，下面讨论其中一些。束搜索（beam search）是一种常用策略，在每一步解码时保留概率最大的 $n$ 个句子，最终选择其中概率最高的一个。

\begin{practicaltip}
束宽通常取 3 至 6 就已足够。继续增大束宽，表现可能反而下降。长度惩罚（length penalty），也称长度归一化（length normalization），是另一项改进，用于补偿束搜索偏好较短句子的倾向。做法是按句子长度调整句子概率，以长度的某个幂作除数。也有方法对生成已用过的词元施加惩罚，以减少重复词元或重复 $n$ 元组。多样化束搜索（diverse beam search）也是一种改进，可以对同一个输入产生多种不同的输出。
\end{practicaltip}

\paragraph*{随机搜索}
基于采样的解码策略以概率方式选择词元，为生成的文本增加多样性与随机性。如果目标是生成丰富且有吸引力的文本，这种策略很有用。例如，给定“I am thirsty. I would like a cup of”作为上下文，下一个词元的概率分布可能偏向“tea”“coffee”或“water”。基于采样的方法仍可能选到概率较低的词，只是机会较小。虽然“tea”的概率最高，“coffee”“water”和“juice”等词仍有机会被选中，使回答更多样。这种方法可用于纯解码器、编码器—解码器和前缀解码器等架构，可灵活适应不同的语言生成任务。

人们对随机采样作了多种改进，减少选中极低概率词的机会，从而改善生成文本的质量。一种改进是温度采样（temperature sampling），计算词表中各词元的概率时，调整 softmax 函数中的温度系数。
\[
P(x_j\mid x_{<i})=\frac{\exp(l_j/t)}{\sum_{j'}\exp(l_{j'}/t)}.
\]
其中 $l'_j$ 表示各词的 logit，$t$ 是温度系数。温度越低，概率较高的词越容易被选中，概率较低的词越不容易被选中。温度为 1 时，默认按原概率分布随机采样。温度趋近 0 时，结果接近贪心搜索。温度无限增大时，采样趋近均匀采样。
\transnote{原书第 59 页的公式用 $l_j$ 表示 logit，紧随其后的文字写作 $l'_j$。}

\begin{practicaltip}
另一项改进是 Top-$k$ 采样（Top-$k$ sampling），截去概率较低的词元，只在概率最高的 $k$ 个词元中采样。Top-$p$ 采样（Top-$p$ sampling），也称核采样（nucleus sampling），是另一种策略。它从累计概率大于或等于指定值 $p$ 的最小词元集合中采样。构造该集合时，按生成概率从高到低加入词元，直到累计概率超过 $p$。例如，按序加入词元，直到累计概率超过 0.8，就只在已加入的词元中采样。
\end{practicaltip}

\subsubsection{三维并行}
三维并行（3D parallelism）把数据并行、流水线并行和张量并行三种主要技术结合起来，以高效训练大语言模型。

\paragraph*{数据并行}
数据并行（data parallelism）把模型参数和训练数据集分布到多个 GPU，以加快训练。每个 GPU 处理分配给自己的数据并计算梯度，再汇总各设备的梯度，用于更新每个 GPU 上的模型。

\begin{practicaltip}
DeepSpeed 程序库引入的 ZeRO 技术解决数据并行中的内存冗余。通常，数据并行使每个 GPU 都要存储同一个大语言模型的完整副本，包括模型参数、梯度和优化器参数（\refcite{ch02-rajbhandari}{Rajbhandari et al., 2020}）。这种重复存储消耗额外内存。ZeRO 让每个 GPU 只保留一部分数据，需要时从其他 GPU 取回其余数据。根据分片存储的内容，提出了三种策略，即优化器状态分片、梯度分片和参数分片。测试表明，前两种不会增加通信开销。第三种使通信开销增加约 50\%，但节省的内存量取决于 GPU 数量。PyTorch 也引入了类似 ZeRO 的全分片数据并行（Fully Sharded Data Parallel，FSDP）技术。
\end{practicaltip}

\paragraph*{流水线并行}
流水线并行（pipeline parallelism）把大语言模型的不同层分配到若干 GPU。为减少数据传输开销，连续的层会分到同一个 GPU。简单的实现可能无法充分利用 GPU。GPipe 和 PipeDream 等进阶方法同时处理多个数据批次并异步更新梯度，以提高效率（\refcite{ch02-harlap}{Harlap et al., 2018}；\refcite{ch02-huang}{Huang et al., 2019}）。

\paragraph*{张量并行}
张量并行（tensor parallelism）将大语言模型的张量或参数矩阵分割，分配到多个 GPU。例如，可在矩阵乘法中按列拆分参数矩阵，交给不同的 GPU 计算，然后合并各 GPU 的结果。Megatron-LM 等程序库支持张量并行，可用于更复杂的张量（\refcite{ch02-shoeybi}{Shoeybi et al., 2019}）。
